% GENERATED FILE. Edit canonical lessons, then run npm run generate:books.
% canonical-source-hash: fnv1a64:c49fd7b0af291fb9
% canonical-lessons: TA-C311-mulakkam, TA-C311-unnaviratam, TA-C311-kilarntelu, TA-C311-elucci, TA-C311-pirakatanam

\chapter{Raising Voices}
\label{ch:ta-raising-voices}
\hlchaptermodality{311}

\begin{chapteropening}
\textbf{By the end of this chapter:} \emph{I can talk about protests and uprisings.}

Complete the last lesson of chapter 311: I can talk about protests and uprisings.
\end{chapteropening}

\section[\texorpdfstring{\ta{முழக்கம்} (muḻakkam)}{முழக்கம் (muḻakkam)}]{\ta{முழக்கம்} (muḻakkam) — a slogan; a roar}

\label{lesson:TA-C311-mulakkam}



\begin{quote}
Before the new one: say the Tamil for to portray, then the Tamil for the media.
\end{quote}

\subsection*{You'll want to know: \ta{முழக்கம்}}
\textbf{\ta{முழக்கம்}} — \emph{muḻakkam} — \textquotedblleft{}a slogan; a roar\textquotedblright{}.

\begin{quote}
\ta{மக்கள்} \ta{தெருவில்} \ta{முழக்கங்களை} \ta{எழுப்பினார்கள்}.
\end{quote}

\begin{quote}
\emph{makkaḷ teruvil muḻakkaṅkaḷai eḻuppiṉārkaḷ.}
\end{quote}

\begin{quote}
People raised slogans in the street.
\end{quote}

\ta{முழக்கம்} \ta{எழுப்பு}, to raise a slogan.

\begin{cousinweb}
\textbf{\ta{முழக்கம்}} (\emph{muḻakkam}) is native Dravidian (DEDR 4989): a loud noise, as of thunder or drums, from \ta{முழங்கு}, \textquotedblleft{}to roar\textquotedblright{}.
\end{cousinweb}

\begin{grammarlens}[title={What you've built}]
One more word for telling what happened.
\end{grammarlens}

\subsection*{Your turn}
\begin{itemize}
\raggedright
  \item \emph{Say it:} \emph{muḻakkam}
  \item \emph{Say it:} \emph{muḻakkam}, once more
  \item \emph{Say it:} \emph{makkaḷ teruvil muḻakkaṅkaḷai eḻuppiṉārkaḷ.} — People raised slogans in the street.
  \item \emph{Recall:} say the Tamil for to portray, then the Tamil for the media, then say \emph{muḻakkam} again
\end{itemize}

\begin{tcolorbox}[breakable,colback=teal!4,colframe=teal!35!black,title={Before you move on}]
What does \ta{முழக்கம்} mean? (\textquotedblleft{}A slogan; a roar\textquotedblright{}.) Say it once more.
\end{tcolorbox}
\section[\texorpdfstring{\ta{உண்ணாவிரதம்} (uṇṇāviratam)}{உண்ணாவிரதம் (uṇṇāviratam)}]{\ta{உண்ணாவிரதம்} (uṇṇāviratam) — a hunger strike; a fast}

\label{lesson:TA-C311-unnaviratam}



\begin{quote}
Before the new one: say the Tamil for the media, then the Tamil for a slogan.
\end{quote}

\subsection*{You'll want to know: \ta{உண்ணாவிரதம்}}
\textbf{\ta{உண்ணாவிரதம்}} — \emph{uṇṇāviratam} — \textquotedblleft{}a hunger strike; a fast\textquotedblright{}.

\begin{quote}
\ta{விவசாயிகள்} \ta{மூன்று} \ta{நாள்} \ta{உண்ணாவிரதம்} \ta{இருந்தார்கள்}.
\end{quote}

\begin{quote}
\emph{vivacāyikaḷ mūṉṟu nāḷ uṇṇāviratam iruntārkaḷ.}
\end{quote}

\begin{quote}
The farmers went on a three-day hunger strike.
\end{quote}

The verb is \ta{இரு}: \ta{உண்ணாவிரதம்} \ta{இரு}, to fast in protest.

\begin{cousinweb}
\textbf{\ta{உண்ணாவிரதம்}} (\emph{uṇṇāviratam}) joins native \ta{உண்ணா}, \textquotedblleft{}not eating\textquotedblright{}, from \ta{உண்}, \textquotedblleft{}to eat\textquotedblright{} (DEDR 600), and \ta{விரதம்}, \textquotedblleft{}a vow\textquotedblright{}, borrowed from Sanskrit \emph{vrata} (CDIAL 12228).
\end{cousinweb}

\begin{grammarlens}[title={What you've built}]
One more word for telling what happened.
\end{grammarlens}

\subsection*{Your turn}
\begin{itemize}
\raggedright
  \item \emph{Say it:} \emph{uṇṇāviratam}
  \item \emph{Say it:} \emph{uṇṇāviratam}, once more
  \item \emph{Say it:} \emph{vivacāyikaḷ mūṉṟu nāḷ uṇṇāviratam iruntārkaḷ.} — The farmers went on a three-day hunger strike.
  \item \emph{Recall:} say the Tamil for the media, then the Tamil for a slogan, then say \emph{uṇṇāviratam} again
\end{itemize}

\begin{tcolorbox}[breakable,colback=teal!4,colframe=teal!35!black,title={Before you move on}]
What does \ta{உண்ணாவிரதம்} mean? (\textquotedblleft{}A hunger strike; a fast\textquotedblright{}.) Say it once more.
\end{tcolorbox}
\section[\texorpdfstring{\ta{கிளர்ந்தெழு} (kiḷarnteḻu)}{கிளர்ந்தெழு (kiḷarnteḻu)}]{\ta{கிளர்ந்தெழு} (kiḷarnteḻu) — to rise up}

\label{lesson:TA-C311-kilarntelu}



\begin{quote}
Before the new one: say the Tamil for a slogan, then the Tamil for a hunger strike.
\end{quote}

\subsection*{You'll want to know: \ta{கிளர்ந்தெழு}}
\textbf{\ta{கிளர்ந்தெழு}} — \emph{kiḷarnteḻu} — \textquotedblleft{}to rise up\textquotedblright{}.

\begin{quote}
\ta{அந்த} \ta{முடிவுக்கு} \ta{எதிராக} \ta{மக்கள்} \ta{கிளர்ந்தெழுந்தார்கள்}.
\end{quote}

\begin{quote}
\emph{anta muṭivukku etirāka makkaḷ kiḷarnteḻuntārkaḷ.}
\end{quote}

\begin{quote}
The people rose up against that decision.
\end{quote}

Past (I …): \ta{கிளர்ந்தெழுந்தேன்}, \emph{kiḷarnteḻuntēṉ}.

A written word: \ta{கிளர்}, to rise up + \ta{எழு}, to rise.

\begin{cousinweb}
\textbf{\ta{கிளர்ந்தெழு}} (\emph{kiḷarnteḻu}) is a Tamil compound built from \ta{கிளர்}, \textquotedblleft{}to rise, emerge, be stirred\textquotedblright{} (DEDR 1583), and \ta{எழு}, \textquotedblleft{}to rise\textquotedblright{} (DEDR 851).
\end{cousinweb}

\begin{grammarlens}[title={What you've built}]
One more word for telling what happened.
\end{grammarlens}

\subsection*{Your turn}
\begin{itemize}
\raggedright
  \item \emph{Say it:} \emph{kiḷarnteḻu}
  \item \emph{Say it:} \emph{kiḷarnteḻu}, once more
  \item \emph{Say it:} \emph{anta muṭivukku etirāka makkaḷ kiḷarnteḻuntārkaḷ.} — The people rose up against that decision.
  \item \emph{Recall:} say the Tamil for a slogan, then the Tamil for a hunger strike, then say \emph{kiḷarnteḻu} again
\end{itemize}

\begin{tcolorbox}[breakable,colback=teal!4,colframe=teal!35!black,title={Before you move on}]
What does \ta{கிளர்ந்தெழு} mean? (\textquotedblleft{}To rise up\textquotedblright{}.) Say it once more.
\end{tcolorbox}
\section[\texorpdfstring{\ta{எழுச்சி} (eḻucci)}{எழுச்சி (eḻucci)}]{\ta{எழுச்சி} (eḻucci) — an uprising; an upsurge}

\label{lesson:TA-C311-elucci}



\begin{quote}
Before the new one: say the Tamil for a hunger strike, then the Tamil for to rise up.
\end{quote}

\subsection*{You'll want to know: \ta{எழுச்சி}}
\textbf{\ta{எழுச்சி}} — \emph{eḻucci} — \textquotedblleft{}an uprising; an upsurge\textquotedblright{}.

\begin{quote}
\ta{நாடு} \ta{முழுவதும்} \ta{ஒரு} \ta{பெரிய} \ta{எழுச்சி} \ta{நடந்தது}.
\end{quote}

\begin{quote}
\emph{nāṭu muḻuvatum oru periya eḻucci naṭantatu.}
\end{quote}

\begin{quote}
A great uprising took place across the country.
\end{quote}

From \ta{எழு}, to rise.

\begin{cousinweb}
\textbf{\ta{எழுச்சி}} (\emph{eḻucci}) is native Dravidian: a noun from \ta{எழு}, \textquotedblleft{}to rise, be aroused\textquotedblright{} (DEDR 851).
\end{cousinweb}

\begin{grammarlens}[title={What you've built}]
One more word for telling what happened.
\end{grammarlens}

\subsection*{Your turn}
\begin{itemize}
\raggedright
  \item \emph{Say it:} \emph{eḻucci}
  \item \emph{Say it:} \emph{eḻucci}, once more
  \item \emph{Say it:} \emph{nāṭu muḻuvatum oru periya eḻucci naṭantatu.} — A great uprising took place across the country.
  \item \emph{Recall:} say the Tamil for a hunger strike, then the Tamil for to rise up, then say \emph{eḻucci} again
\end{itemize}

\begin{tcolorbox}[breakable,colback=teal!4,colframe=teal!35!black,title={Before you move on}]
What does \ta{எழுச்சி} mean? (\textquotedblleft{}An uprising; an upsurge\textquotedblright{}.) Say it once more.
\end{tcolorbox}
\section[\texorpdfstring{\ta{பிரகடனம்} (pirakaṭaṉam)}{பிரகடனம் (pirakaṭaṉam)}]{\ta{பிரகடனம்} (pirakaṭaṉam) — a proclamation}

\label{lesson:TA-C311-pirakatanam}



\begin{quote}
Before the new one: say the Tamil for to rise up, then the Tamil for an uprising.
\end{quote}

\subsection*{You'll want to know: \ta{பிரகடனம்}}
\textbf{\ta{பிரகடனம்}} — \emph{pirakaṭaṉam} — \textquotedblleft{}a proclamation\textquotedblright{}.

\begin{quote}
\ta{அரசு} \ta{ஒரு} \ta{முக்கியமான} \ta{பிரகடனம்} \ta{வெளியிட்டது}.
\end{quote}

\begin{quote}
\emph{aracu oru mukkiyamāṉa pirakaṭaṉam veḷiyiṭṭatu.}
\end{quote}

\begin{quote}
The government issued an important proclamation.
\end{quote}

The verbs are \ta{செய்} or \ta{வெளியிடு}: \ta{பிரகடனம்} \ta{செய்}.

\begin{cousinweb}
\textbf{\ta{பிரகடனம்}} (\emph{pirakaṭaṉam}) is borrowed from Sanskrit \emph{prakaṭana}, \textquotedblleft{}manifesting, bringing to light\textquotedblright{} (Monier-Williams), from \emph{prakaṭa}, \textquotedblleft{}manifest\textquotedblright{} (CDIAL 8428).
\end{cousinweb}

\begin{grammarlens}[title={What you've built}]
One more word for telling what happened.
\end{grammarlens}

\subsection*{Your turn}
\begin{itemize}
\raggedright
  \item \emph{Say it:} \emph{pirakaṭaṉam}
  \item \emph{Say it:} \emph{pirakaṭaṉam}, once more
  \item \emph{Say it:} \emph{aracu oru mukkiyamāṉa pirakaṭaṉam veḷiyiṭṭatu.} — The government issued an important proclamation.
  \item \emph{Recall:} say the Tamil for to rise up, then the Tamil for an uprising, then say \emph{pirakaṭaṉam} again
\end{itemize}

\begin{tcolorbox}[breakable,colback=teal!4,colframe=teal!35!black,title={Before you move on}]
What does \ta{பிரகடனம்} mean? (\textquotedblleft{}A proclamation\textquotedblright{}.) Say it once more.
\end{tcolorbox}
